\documentclass[aspectratio=169]{beamer}

\usetheme{Warsaw}
\usecolortheme{whale}

\setbeamertemplate{headline}{}
\setbeamertemplate{navigation symbols}{}
\setbeamercolor{background canvas}{bg=}

\usepackage{beamerthemesplit}
\usepackage{pgf}
\usepackage{epsfig}
\usepackage{epstopdf}
\usepackage{graphicx}
\usepackage{url}
\usepackage{listings}
\usepackage{amsmath}

\usepackage[utf8]{inputenc}
\usepackage{polski}

\graphicspath{ {./gfx/} }

\author[Amadeusz Kosik]{Amadeusz Kosik}
\title[Interwalled \insertframenumber/\inserttotalframenumber]{Implementacja efektywnego obliczania złączeń interwałowych \\ dla biblioteki Apache Spark}

% Maybe useful:
% \begin{frame}[allowframebreaks,fragile]
% \begin{verbatim} ... \end{verbatim}
% \item [$+$] 
% \item [$-$] 


% Wersja 1: cel pracy, wprowadzenie do tematyki, przegląd istniejących rozwiązań,
% uzasadnienie potrzeby pracy; analiza zagadnienia, teoria (jeśli dotyczy), trudności
% związane z zadaniem, idea rozwiązania

\date{Warszawa, 2026}
\institute{Promotor pracy: \\ dr hab. inż. Tomasz Gambin, profesor uczelni}

\begin{document}

\frame{\titlepage}

\frame{\tableofcontents\frametitle{Plan prezentacji}}

\section{Definicja zadania}
	\begin{frame} 
		\frametitle{Definicja zadania: słowna}
		\parbox{\linewidth}{Zadana jest aplikacja, która dla dwóch zbiorów $A$ i $B$ interwałów zwróci zbiór wszystkich par takich, gdzie nakładających 
		się interwałów takich, gdzie lewy element pary pochodzi z $A$, prawy pochodzi z $B$ oraz interwały w parze nakładają się na siebie.}
		\vfill
		\parbox{\linewidth}{Aplikacja ma działać na środowisku Apache Spark, aby możliwe było jej wpięcie w~bardziej złożone potoki przetwarzania danych
		zbudowane na bazie tejże biblioteki.}
		\vfill
		\parbox{\linewidth}{Oczekiwana jest lepsza wydajność lub większa skalowalność w porównaniu do już dostępnych rozwiązań.}
	\end{frame}
	
	\begin{frame} 
		\frametitle{Definicja zadania: precyzyjna}
		Dla interwału zdefiniowanego jako $a = (a_{od}, a_{do}, a_1, ..., a_k); a_{od} \leq a_{do}$, gdzie $a_{od}$ i $a_{do}$ są danymi liczbowymi, 
		zbiorów interwałów $A$ i $B$ należy zwrócić zbiór par interwałów taki, że: 
		\begin{equation*}
			 f(A, B) = \left\lbrace \left(\left(a_{od}, a_{do}, ... a_n\right), \left(b_{od}, b_{do}, ... b_n\right)\right):
			 a \in A \land b \in B \land \left( a_{od} \leq b_{do} \land a_{do} \geq b_{od} \right) \right\rbrace
		\end{equation*}
	\end{frame}
	
	\begin{frame} 
		\begin{table}
			\begin{minipage}{.45\textwidth} \centering
				\begin{tabular}{r|rr}
				      & $a_{od}$ & $a_{do}$ \\ \hline
				$a_1$ &        2 &        6 \\
				$a_2$ &        9 &       12 \\
				$a_3$ &       10 &       16 \\
				$a_4$ &       24 &       27 \\ \hline \hline
				      & $b_{od}$ & $b_{do}$ \\ \hline
				$b_1$ &        2 &       10 \\
				$b_2$ &       14 &       16 \\
				$b_3$ &       18 &       20 \\
				$b_4$ &       27 &       30 
				\end{tabular}	
			\end{minipage}%		
			\hspace{0.05\textwidth}
			\begin{minipage}{.45\textwidth} \centering
				\begin{tabular}{r|rr||r|rr}
				      & $a_{od}$ & $a_{do}$ &       & $b_{od}$ & $b_{do}$ \\ \hline
				$a_1$ &        2 &        6 & $b_1$ &        2 &       10 \\
				$a_2$ &        9 &       12 & $b_1$ &        2 &       10 \\
				$a_3$ &       10 &       16 & $b_1$ &        2 &       10 \\
				$a_3$ &       10 &       16 & $b_1$ &        2 &       10 \\
				$a_3$ &       10 &       16 & $b_2$ &       14 &       16 \\
				$a_4$ &       24 &       27 & $b_4$ &       27 &       30 			
				\end{tabular}
			\end{minipage}	
			\caption{Przykładowe dane wejściowe dla aplikacji i~oczekiwany dla nich rezultat.}
		\end{table}	
	\end{frame}


\section{Apache Spark}
	\begin{frame}
		\frametitle{Apache Spark}
		TUTAJ OPIS
	\end{frame}

\section{Istniejące aplikacje}
	\begin{frame}
		\frametitle{Sequila}
		TUTAJ OPIS
	\end{frame}
	
	\begin{frame}
		\frametitle{???}
		POSZUKAĆ WIĘCEJ
	\end{frame}
	
\section{Założenia rozwiązania}
	\begin{frame}
		\frametitle{Augumented Interval List}
		TUTAJ OPIS ALGORYTMU
	\end{frame}
	
	\begin{frame}
		\frametitle{Augumented Interval List}
		TUTAJ DLACZEGO AKURAT TEN
	\end{frame}

\section{Weryfikacja}
	
	\begin{frame}
		\frametitle{Weryfikacja poprawności}
		Testy syntetyczne
	\end{frame}
	
	\begin{frame}
		\frametitle{Weryfikacja poprawności}
		Testy vs inne biblioteki
	\end{frame}

	\begin{frame}
		\frametitle{Weryfikacja wydajności}
		Punkt odniesienia: Sequila
	\end{frame}

	\begin{frame}
		\frametitle{Weryfikacja wydajności}
		Punkt odniesienia: BioDataGeeks
	\end{frame}

\end{document}
